\documentclass[12pt, a4paper]{article}

\usepackage{ctex}
\usepackage{float, booktabs, parskip}
\usepackage{amsmath, amsthm, amssymb, mathrsfs, bm}
\usepackage{geometry}
\usepackage{hyperref}

\geometry{a4paper, left=2cm, right=2cm, top=2.5cm, bottom=2.5cm}

\setlength{\parindent}{2em}
\setlength{\parskip}{0.5em}
\allowdisplaybreaks[4]



\title{{\large{\textbf{极地系统的观测、认知与战略——《极地环境快变与突变遥感》讲座报告}}}}
\author{刘滨瑞 \\ 地球系统科学系 \quad \href{mailto:lbr21@mails.tsinghua.edu.cn}{lbr21@mails.tsinghua.edu.cn}}
\date{\today}

\begin{document}

\maketitle



本报告内容基于程晓老师10月16日的讲座《极地环境快变与突变遥感》。
讲座呈现了中山大学极地科考团队在南极DNT系统的观测成果和在北极建立的“空-天-冰-海”一体化观测体系，揭示了极地系统在全球变暖背景下正经历的深刻变化。
另外，程老师也从自身经历出发，介绍了中国极地科考的情况。



\section{Drygalski冰舌: 如何理解亚稳态与临界点}

程晓老师团队详细重构了2022年的\textbf{Drygalski冰舌崩解事件}。
2022年1月15日，汤加火山剧烈爆发，向平流层注入了前所未有的水汽量。
这个发生在南太平洋的地质事件，通过大气和海洋波动，在数千公里外的南极引发了一系列海啸。
1月16日，海啸抵达Drygalski冰舌。遥感影像和验潮站的数据能够精确地相互印证: 冰舌前端的应变响应与海啸的到达时刻完美吻合。

程老师指出，\textbf{冰舌前端在崩解之前就已经存在一条几乎横贯整个冰舌的裂缝}。这条裂缝不是海啸造成的，而是冰舌长期演化过程中积累的结构性脆弱点。
海啸到达后，裂缝在两小时内迅速扩张，增长超过1km，最终导致冰舌前端一块约10kmx5km的冰体崩解下来。
这一洞察揭示了一个深刻的科学问题: \textbf{海啸本身只是冰舌崩解的诱因，根本原因在于冰舌在长期的应力作用下已经进入了一个高度不稳定的状态}。
即使没有汤加火山爆发，也许下一次强风暴、下一次潮汐异常、甚至下一次季节性的温度波动，都可能触发同样的崩解。

这引发了我的思考。
我们对“气候临界点”的传统认知可能过于简化。
我们习惯于将临界点理解为一个明确的阈值——温度升高多少度、海冰减少多少百分比——一旦跨越就会发生突变。
但Drygalski冰舌的案例告诉我们，“临界点”也可能是一个过程，意味着进入了某种\textbf{亚稳态 (metastable state)}。
此时系统实际上已经失去了韧性，任何一个扰动都可能成为“压垮骆驼的最后一根稻草”。
问题在于，我们能否识别出一个系统已经进入了这种亚稳态？

对于冰架崩解这一案例而言。
程老师团队通过冰流数值模式模拟了崩解事件对该冰架的动力学影响。
另外，程老师介绍的冰架崩解遥感监测技术——结合合成孔径雷达 (SAR)、光学卫星、验潮站等多源数据——在事后归因方面已经相当成熟。
借助这些数据和方法，我们应该如何在崩解之前识别出“隐藏的裂缝”，从而做出预警？
这个问题的核心困难在于，冰体内部的应力积累过程是不可见的。
卫星遥感只能捕捉冰舌表面的变形。
当裂缝扩展到表面时，我们已经错过了预警的最佳时机。
那条横贯Drygalski冰舌的裂缝，很可能在冰体深部存在了数年甚至数十年，在应力逐年累积的过程中缓慢扩展，直到海啸的到达。

从复杂系统 (Complex System) 的角度，亚稳态系统通常会展现出某些可探测的早期预警信号。
例如，冰舌对潮汐、风暴等外部扰动的响应方式可能发生微妙变化，例如恢复时间变长、波动幅度增大、空间变形模式出现异常等。
程老师展示的验潮站数据与遥感影像已经足以捕捉到部分的系统信号了，但是仍然需要更进一步的发展。
在观测方法方面，我们需要更先进的冰体内部探测技术。例如，我们可以通过地震波探测捕捉微裂缝扩展时释放的声发射信号；或者通过分析冰体对海洋潮汐的非线性响应，来反演内部损伤的分布。
在理论方面，我们还需要建立从微观损伤到宏观失稳的跨尺度理论，从而解释一条毫米级的微裂缝如何演化成公里级的崩解的。
而在数据方面，我们应该将有限的观测数据与高保真数值模拟深度融合，通过数据同化技术不断更新对冰体状态的评估。



\section{DNT系统: 耦合过程与关键科学问题}

讲座中，程晓老师详细介绍了团队正在开展的DNT系统研究。
DNT即Drygalski冰舌-Nansen冰架-Terra Nova Bay冰间湖系统，这是南极西侧一个空间尺度约数百公里的区域。
DNT系统在一个相对紧凑的空间内集中展现了极地系统中几乎所有关键的耦合过程 (大气-冰盖-海洋)。
Drygalski冰舌作为一个长达70km的固定冰结构，对其下游的冰架和海洋环流起到了屏障作用。
冰舌的存在改变了局地海洋环流模式，阻挡了暖水的入侵，从而保护了其后方的Nansen冰架免受底部融化。
Terra Nova Bay冰间湖是南极最活跃的海冰生产区之一。
在冬季，强烈的下坡风驱动海冰向外输送，持续暴露开阔水面。
海冰的快速生成释放出盐分，形成高盐、高密度的陆架水，这是南极底层水 (AABW) 的重要源区。
Nansen冰架位于Drygalski冰舌和Terra Nova Bay冰间湖之间，其底部融化速率受到两侧海洋环流的夹击，同时又受到地表下坡风强度的调控。

程晓老师提出的核心科学问题非常明确且具有挑战性:
\begin{enumerate}
    \item 为何地处75°S高纬度的冰架会发育出成熟的表面与内部水文系统？
    \item 固定冰支撑与冰舌崩解机制研究
    \item 冰间湖对冰架、冰舌崩解的响应机制研究
    \item 冰架、冰舌崩塌的建模与预测
\end{enumerate}



\section{冰盖融水: 尺度鸿沟与观测困境}

极地研究的一个重要问题是: 冰架如何失去质量的？
程晓老师在讲座中指出，我们目前面临着两个核心难题: 1. \textbf{表面融化定量化不足}；2. \textbf{底部过程机制不清}。
其中的困难在于: \textbf{关键物理过程发生的尺度，与我们观测和模拟能力覆盖的尺度之间存在着巨大的鸿沟}。

程晓老师展示的卫星观测数据显示，南极冰盖表面融水生产正在快速增加。
程老师团队发展的1992-2022/2023年格陵兰与南极冰盖逐日融水体积数据集，是该领域研究的重要进展。
但程老师同时指出，我们还需要更精确地反演融化量，才能评估融水对冰架稳定性的影响。
表面融化是一个跨越多个尺度的复杂过程。
在宏观尺度 (数公里到数百公里)，我们可以用能量平衡方程估算总融化量。
但是，融水一旦产生，它的去向难以被解析。
例如，渗入积雪被重新冻结、沿冰面流动形成融水河、注入裂隙触发水力压裂这些关键过程，它们均发生在米到十米的尺度上。
当前卫星遥感的空间分辨率和数值模式的网格尺度，都难以直接刻画这些过程。

更关键的是，这些小尺度融水过程对大尺度的冰架稳定性具有阈值效应。
一个宽度仅几米的裂隙，一旦充满融水，在特定应力条件下可能在几小时内扩展到公里尺度，触发整个冰架的崩解。
2002年Larsen B冰架的崩解正是这一机制的典型案例。
这意味着，即使我们准确估算了融化总量，如果不知道融水的空间分布和动态演化路径，仍然无法预测崩解事件的发生。

与表面融化相比，底部融化过程更加隐蔽，却是质量损失的主导项。
程晓老师提出，对于南极和格陵兰的大部分冰架而言，底部融化的贡献往往数倍于表面融化和冰山崩解的总和。
但是，我们对其机制仍然知之甚少。
海洋如何融化冰架底部？暖水如何穿过密度锋进入冰架下方的空腔？融化速率如何随深度、距离和海洋温度分布变化？融化释放的淡水又如何反馈到局地海洋分层，进而影响融化速率本身？
这一系列问题的答案，我们尚不知道。

冰架下方是一个极难观测的环境。
程晓老师向我们展示的“悟空6000”无人潜水器实现了我国首次冰架底部综合调查，这是技术上的重大突破。
但即使是如此先进的设备，也只能在有限的时间窗口内获取若干离散测点的数据。
考虑到冰架底部的融化过程在空间上高度不均匀，在时间上也有强烈的季节性和年际变化。
因此，现有的观测数据是远远不足以支撑我们对冰架底部融化过程的认知的。

我认为，无论是表面融化还是底部融化，我们势必需要将观测与模拟深度融合。
程晓老师展示的无人机遥感、冰雷达等技术正在填补部分的尺度鸿沟，但同时我们也要发展新的理论框架，即如何从有限的、离散的观测数据中，通过数据同化技术重构整个融化过程的时空演化，以及如何在物理模型中合理表征那些发生在次网格尺度的关键过程。
在这一方面，我相信人工智能技术或许能够提供新的思路。



\section{从追随到引领: 中国极地战略的演进}

程晓老师向我们展示的\textbf{中山大学“极地号”}科考船，是我国第三艘具备极地破冰能力的科考船，也是第一艘由大学自主运营的极地科考船。
该船总长78.85米、宽17.22米、型深9.7米、满载排水量5852吨、破冰能力为PC4级。
我认为，“极地号”的意义不仅在于技术参数本身，更在于它标志着中国极地科考从“借船出海”到“自主航行”的根本转变。
在“极地号”之前，中国极地科考主要依赖自然资源部中国极地研究中心运营的“雪龙号”和“雪龙2号”，航次计划需统筹考虑站点补给、多学科需求，科学家难以获得足够船时来实施长期、连续、针对性的观测。
而中山大学拥有自己的科考船，意味着科研团队可以根据科学问题灵活设计航次、布置观测网络、进行适应性调整。
程老师团队2024年的北极科考任务——8月17日进入北极圈、9月14日离开北极圈、在北极海域工作29天、冰区航程1473海里、最高纬度85°N——正是这种自主能力的体现。

与“极地”号的水面突破相呼应，国产无人潜水器\textbf{“悟空6000”}成功完成了在北极85°N、水深6000米的冰下作业，实现了我国首次冰架底部综合调查。
“悟空6000”在冰下作业时没有GPS、没有实时通讯，完全依靠惯性导航和声学定位自主完成观测任务并准确返回，这对控制算法、能源管理、传感器可靠性都提出了极高要求。
冰上与冰下观测能力的同步突破，标志着中国在极地科考硬件上已跻身世界前列。

但技术能力只是第一步。如果说“极地”号和“悟空6000”代表了硬实力，那么\textbf{FACE (Following Arctic/Antarctic ICE) 国际计划}则代表了软实力，即科学议程的设置权。
程老师介绍到，FACE计划于2023年启动，由中国科学家主导建立科学委员会和科学顾问小组，成员包括来自挪威、中国、加拿大、冰岛、德国等多国的极地科学家。
FACE是一个长期、系统、多国参与的科学计划，其三阶段航次规划包括: 2024-2025年关注大西洋扇区、2026-2028年扩展到太平洋扇区、2029-2032年拓展至北大西洋。
在极地科学领域，谁能发起和主导这样的国际大科学计划，谁就掌握了话语权，谁就能决定“什么是重要的科学问题”、“观测资源如何配置”、“数据如何共享”。
这可能比单纯的技术能力更加重要。

FACE具有一个跨圈层、多尺度、过程导向的前瞻性科学框架，包括大气(高空大气、极端天气、北极云和降水)、冰圈(边缘冰区、海冰输出、融池、冰盖)、海洋(水团、中尺度涡、碳循环)、环境与生态(微塑料、持久性污染物、水下噪声、微生物多样性)。
其核心逻辑是在全球变暖背景下，通过天-空-冰-海潜协同立体观测，“追踪”不断消失的极地海冰和快速崩解的冰架，理解其变化机理，预测其未来突变。

与国际科研合作相对的，是单边主义、保护主义和地缘政治博弈。
美国、加拿大、挪威、冰岛、德国等国，都纷纷推出了自己的极地科考计划，试图在极地科学领域占据主导地位。
这种竞争的激烈程度，从程晓老师在“极地号”上的科考经历之中可窥一斑。
美国、加拿大媒体多次报道“极地号”进入北冰洋；
在"极地号"穿越对马海峡和白令海峡的时候，遭到了美国、加拿大、日本、俄罗斯的巡逻船和巡逻飞机的跟踪与监视。
这反映了极地科考背后复杂的现实利益博弈。

对于中国而言，北极具有重要的科学、经济和地缘政治意义。
从科学层面，中国作为最大的碳排放国和人口大国，理解极地系统对全球气候的影响，是制定科学政策的基础。
从经济层面，北极航道开通将使中国到欧洲的海运距离缩短约30\%，这关乎能源安全与供应链多元化。
但航道利用需要更精准的海冰预报，这需要长期、系统的科学观测作为支撑。
从地缘政治层面，北极是少数几个尚未完全被主权瓜分的区域，《斯瓦尔巴条约》赋予签约国的科学考察权利在这里将直接转化为国际话语权。
面对这一复杂局面，中国极地科考呈现出一种看似矛盾的双重选择：一方面，通过发起FACE这样的开放性国际合作计划，强调科学研究的公共产品属性，展现负责任大国的姿态；另一方面，通过自主建造“极地”号、发展“悟空6000”等关键技术，确保在核心能力上不受制于人。
这既是科学与政治的微妙平衡，也是战略视野拓展所带来的必然选择。

\end{document}
